% 15_ai_method
\section{AI 辅助开发方法与正确性保障}

本队采用一套以机器可判定的判定门为核心的 AI 辅助开发方法。方法把开发分工划为两侧：团队一侧确定架构方向、划定原子改动的边界、制定验收标准；AI 一侧在真机与 QEMU 的闭环内反复迭代实现，把每一次改动运行至判定门，由判定门给出通过与否的判定。其关键在于全部验收标准都落实为机器可判定的判定门，一次运行即产出一个可比对的数值，以该数值取代肉眼读取日志得出结论的环节。以往判断一次内核改动是否正确，依赖人工阅读串口日志逐条比对，既慢又易漏；现将正确性的判据固化为数值化的判定门之后，判定既快也可复核。

\subsection{方法的架构}

这套方法由三个部分构成：团队掌握的规格与验收、AI 在闭环内的实现迭代，以及横在两者之间的一组机器可判定的判定门。团队给出一次改动的目标、边界与验收阈值；AI 将实现运行至判定门，读取判定门返回的布尔判定与数值，据此继续迭代或收敛；判定门运行在板上或 QEMU 上，产出的数值同时用于放行改动与否定设想。图~\ref{fig:aimethod} 给出这一双侧分工与中间的四类判定门。

\figphl{../figures/out/Faimethod.png}{团队掌握规格与验收、AI 在真机闭环内迭代实现，中间由四类机器可判定的判定门裁决}{fig:aimethod}

判定门指一段能够自行判定对错的运行：给定输入，运行完整流程，将结果与同一条基线比对，输出一个表示通过与否的布尔量以及一个可追溯的数值。sysbench 与 schbench 是既有的上游基准，Linux 同机数据是既有的比对基线；将其组织为一次运行即给出放行判据的判定门，以及贯穿全部改动的评审与暖重启回路，是本队新增的部分。

\subsection{四类判定门}

第一类是 sysbench 对等回归门。每一次触及内存或调度的改动，都在板上重跑 cpu、threads、mutex、memory 的完整矩阵，与 Linux 同机数据逐轴求比值，任一轴掉出既定区间即判为失败。DVFS 的四步 OPP 提升依靠这道门逐步放行，每一步以 sysbench cpu 8 线程无压降为通过条件，最终 8 线程达到 \keyres{\nCpuTEight 事件/秒，为 Linux \nCpuTEightLin 的 \nCpuTEightRatio}。

第二类是 schbench 延迟阈值门。唤醒延迟相关的改动完成后，以 schbench 测量 p50/p99 唤醒延迟，直接与一个数值阈值比对。跨核唤醒发送 IPI 的原有路径把 p50 抬到约 1042 微秒，wake\_affine 改动完成后降至 \keyres{\nWake 微秒，接近 Linux 的 \nWakeLin 微秒}，该数值即为放行凭据。

第三类是板上无死锁的多智能体对抗式评审门。触及 SMP 同步的改动落地之前，先拆为单一关注点的原子改动，再交由多个相互独立的评审智能体对抗式复核，每次改动配备 7 到 26 名评审，正是这道门查出了两处高危缺陷。其一，DVFS 的 apply\_opp 将一次电压写入失败静默吞没，可能在提升频率之后使核心停留于高频欠压这一存在硬件风险的状态，\ours{改为失败即中止}之后，出错至多导致过压；其二，perf 组采样读路径的一处 use-after-free 在复核中被定位并修复。这类缺陷多数潜伏于 ArceOS 家族共享的跨内核 crate，能够通过随手测试而在 SMP 叠加 COW、THP 时损坏内存，页表与 TLB 一类的修复同时惠及全部三个系统。

第四类是 s\_frac 计时正确性 oracle。CPU 时间记账既要快也要准。这道门以 rusage 的 s\_frac，即系统态时间占比，作为可比对的正确性判据。一处遗漏的 retick 曾把被调度出去的睡眠时段计入 utime/stime，整秒乃至整小时的睡眠都被记为 CPU 时间；以一个无锁的 resume\_ns 下界修复之后，s\_frac 落在 \keyres{0.71，接近 Linux 的 0.72}，表明记账既准确又高效。

\subsection{真机闭环与带外电源控制}

这些判定门要产生价值，前提是能在真实 RK3588 上低成本地反复运行，且这条闭环须无人值守。boardctl.py \ours{把常规一轮迭代收敛为暖重启}：从 StarryOS shell 执行 reboot -f 返回 Linux，经 ssh 触发重启进入 U-Boot，以 fatload 载入部署的映像并 bootm，配合部署、同步、校验与实时失败监视，把一轮迭代压缩到秒级且无需断电。暖重启承担常规路径，但它有一处结构性盲区：只在被测系统仍能响应软件命令时才能恢复。这条闭环恰恰用于验证会令内核挂死的改动，SMP 死锁、早期 panic 与活锁都在其列；一旦某次改动在 shell、SSH 与 U-Boot 三者均不可达之前就挂死，暖重启无从恢复，回路只能停下等人到场物理断电。一套用于验证挂死的闭环，其恢复手段若本身依赖被测系统不挂死，便不成其为无人值守。

\figphl{../figures/out/Fharness.png}{带外电源控制的真机自治测试台：主机与串口适配器接常供电，板卡专用电源经计量式智能插座切换；暖重启为常规快路径，智能插座硬断电冷启为与被测软件状态无关的恢复原语，并解锁上电 TTFI 与墙上功率计量}{fig:harness}

要让闭环真正无人值守，恢复手段必须独立于被测系统的软件状态。板级唯一具备这一性质的原语是电源。为此在板卡供电与市电之间接入\ours{一只可编程的计量式智能插座（米家智能插座）}，主机侧的回路经其 miIO 本地协议在任意状态下对板卡冷复位，全程不经云端。恢复据此分四级逐级退化：L0 暖重启，被测系统存活，秒级，承担常规一轮；L1 智能插座硬断电冷启，被测系统挂死而主机存活，与被测软件状态无关，可从任何挂死恢复；L2 经 U-Boot 启动计数回退到金身映像，用于部署的内核连冷启都过不去时；L3 上报人工，仅当金身映像冷启仍失败，此类情形属真实硬件故障，是这条回路唯一需要人的环节。

图~\ref{fig:harness} 给出这套测试台的拓扑，它令暖重启的两处顾虑一并消解。插座只切板卡专用电源一路，主机与 USB 串口适配器另接常供电。断板电因此不会令主机侧串口适配器重新枚举，控制台日志在复位间保持连续，即便某次改动在 SSH 建立之前就挂死，回路仍能读到确切的挂死点；专用且足额的电源为 8 核冷启提供浪涌余量。硬断电只在少数挂死时触发，常规仍走暖重启，反复上下电既慢又损板的代价由此不再构成约束。

该插座同时带功率计量，为闭环解锁两项此前缺失的量度。其一，TTFI 的口径是上电到首帧推理，唯有一次真实上电事件方能诚实测得；插座使回路能自市电接通计时至首帧推理，给出口径自上电起算的 TTFI，不再以 U-Boot 交接为起点。其二，与 sysbench 及推理负载同步采样插座的墙上功率，即可给出每焦耳事件数对 Linux 的能效比对，正是可复现一节列为最高价值的待补测轴。上述机制现已就位，两项数值随整机复测补齐。四类判定门中凡标注为板上的，都运行在这条测试台上，QEMU 先行验证逻辑，真机负责给出最终数值。

\subsection{测量驱动的对抗式验证否定了若干设想}

判定门给出的数值既用于放行改动，也用于否定基于推断的设想，后者同样重要。DDR 带宽差距最初被推测为 THP 或 tick 记账所致，经板上 schbench 与 sysbench 测量，两个假设都被数据否定，真正的根因是 DDR 总线固定于上电频率、变频由 Rockchip 的 SIP DRAM 接口负责而主线固件缺少这一 EL3 处理程序，属固件层面的外部因素。fork 在约 250 个进程处出现 EFAULT，曾被判断为缺页处理的问题，板上复现后定位到 COW 引用计数 u8 溢出，\ours{将 u8 改为 u32}即解决。自适应 poll-idle 一度被寄望于胜过固定 50 微秒空转，板上测量表明其并无优势，因而未予采纳。每一个被否定的设想背后都对应一次真机测量。

\subsection{团队掌控架构与验收}

这套方法的边界清晰：AI 负责在闭环内把实现迭代至通过判定门，架构决策与最终验收始终由团队负责。划定原子改动的边界、确定每道判定门的阈值、判定一个数值是否足以放行，均为团队做出的判断。团队成员能在无 AI 在场时口头阐明每一处关键设计的动机，并当场修改对应的内核代码：PVTPLL 环长切换为何能在同一电压下换取更多主频、break-before-make 为何需先执行 dsb ishst 再执行 tlbi、s\_frac 为何可作正确性 oracle，均非黑箱。收尾阶段一轮投机实现排查覆盖了整套出厂改动，未查出投机取巧的实现，剩余唯一的硬性障碍是固件受阻的 DDR 变频这一外部因素。
